\section{Expected outcomes and impact}\label{sec:outcomes}

\paragraph{For the group} Either a validated, low-cost way to stop the degradation of \SI{369.5}{nm}
fibre links (printed cells with oxygen absorbers or a clean purge), or a clear identification of the
cause and a recommendation for fibre or end-face engineering. Both reduce down time, repeated polishing,
and slow power and pointing drifts in the experiments. The same cells can protect fibres at
\SI{399}{nm} and other near-UV wavelengths used in the group.

\paragraph{For the field} Fibre delivery of near-UV light is a widespread practical problem in
trapped-ion and neutral-atom experiments and in UV photonics, but the role of the atmosphere at the end
face has, to our knowledge, not been isolated systematically. A controlled comparison of \Oxy, water,
contaminants and particles at CW near-UV intensities of order \SI{e5}{W/cm^2} is suitable for a short paper
(e.g.\ in \emph{Optics Express}, \emph{Applied Optics}, \emph{Review of Scientific Instruments} or
\emph{EPJ Techniques and Instrumentation}). Open CAD files and procedures make the result easy to adopt.

\paragraph{For the student} Training in experimental design for slow, drift-prone measurements, fibre
optics and UV laser handling, gas handling, CAD and 3D printing for functional parts, data acquisition
and statistics, and surface analysis. This mix of skills is valued in quantum-technology labs and in
industry.
